\capitulocapa{images/cap-7.png}{Como criar o primeiro store?}
\label{chap:primeiro-store}

\section{Um modal que começou em uma página}

\begin{historia}
A equipe adicionou à página de Usuários um botão ``Convidar pessoa''. Ao clicar,
um modal era aberto por um \codigo{useState} local. A implementação era simples e
adequada ao pedido daquele momento.

Na semana seguinte, o mesmo convite passou a ser iniciado pelo Header vazio, por
um atalho na tela inicial e por uma ação rápida da busca. Elevar o estado faria
componentes intermediários transportarem propriedades que não utilizavam.

Alguém propôs criar uma store com tudo que o modal poderia precisar: campos do
formulário, mensagens, permissões e resposta da API. Antes de escrever código, a
equipe fez uma pergunta menor: qual é o menor contrato capaz de coordenar a
abertura e o fechamento?
\end{historia}

\ilustracaopendente{7-1-modal-varias-origens}{Cena do DevPanel com o modal de
convite no centro. Três caminhos apontam para ele: o botão da página de Usuários,
uma ação rápida no Dashboard e a busca do Header. Nenhum caminho deve atravessar
componentes intermediários; todos se conectam à mesma pequena store desenhada ao
lado do modal.}

O primeiro store útil não precisa representar uma área inteira do produto. Ele
precisa possuir uma responsabilidade clara, estado suficiente e ações que
expressem as interações permitidas.

\section{Comece pelo contrato}

Antes de chamar \codigo{create}, escrevemos as perguntas que o restante da
aplicação precisa responder:

\begin{itemize}
  \item o modal está aberto?
  \item qual e-mail pode aparecer preenchido ao abrir?
  \item como iniciar um novo convite?
  \item como encerrar o fluxo?
\end{itemize}

O contrato TypeScript nasce dessas perguntas:

\begin{lstlisting}[caption={Estado e ações do fluxo de convite.}]
type InviteState = {
  isOpen: boolean
  initialEmail: string | null
}

type InviteActions = {
  openInvite: (email?: string) => void
  closeInvite: () => void
}

type InviteStore = InviteState & InviteActions
\end{lstlisting}

Separar os tipos de estado e ações torna a intenção visível. Não é uma exigência
do Zustand, mas ajuda o time a discutir o contrato sem se perder na
implementação.

O formulário digitado pela pessoa pode continuar local ao modal. A resposta da
API continua sendo responsabilidade da ferramenta de consultas. A store guarda
somente o que coordena esse fluxo entre regiões distantes da interface.

\begin{decisao}
Modele o comportamento necessário agora. Campos adicionados ``para o futuro''
aumentam o contrato público e tornam mais difícil descobrir a responsabilidade
real da store.
\end{decisao}

\section{Criando a store}

Com o contrato definido, usamos \codigo{create}:

\begin{lstlisting}[caption={Primeira store completa do fluxo de convite.},label={lst:invite-store}]
import { create } from 'zustand'

export const useInviteStore = create<InviteStore>()((set) => ({
  isOpen: false,
  initialEmail: null,

  openInvite: (email) =>
    set({
      isOpen: true,
      initialEmail: email ?? null,
    }),

  closeInvite: () =>
    set({
      isOpen: false,
      initialEmail: null,
    }),
}))
\end{lstlisting}

A forma \codigo{create<InviteStore>()(...)} informa ao TypeScript o contrato
completo. Se esquecermos uma ação, retornarmos um tipo incorreto ou chamarmos
\codigo{openInvite} com um número, o erro aparece durante o desenvolvimento.

A função criadora recebe \codigo{set}. Ela devolve o estado inicial e as ações.
Quando uma ação chama \codigo{set} com um objeto, Zustand combina essa parcela no
nível superior da store. Por isso não repetimos \codigo{isOpen} ao atualizar
apenas outro campo.

Fechar também limpa \codigo{initialEmail}. Essa pequena regra evita que um
convite futuro herde o destinatário anterior. Quem chama a ação não precisa
lembrar de executar duas atualizações separadas.

\section{Ações devem expressar intenção}

Poderíamos expor uma ação genérica chamada \codigo{setOpen}. Ela receberia um
booleano para abrir ou fechar, mas obrigaria os consumidores a conhecer detalhes
adicionais: quando limpar o e-mail? Como abrir já preenchido? O que significa
chamá-la com \codigo{true} duas vezes?

\codigo{openInvite} e \codigo{closeInvite} descrevem eventos do produto. A ação
pode preservar invariantes enquanto a interface apenas comunica a intenção.

Isso não significa criar uma função para qualquer atribuição. Uma ação merece
existir quando dá nome a um comportamento, reúne campos que precisam mudar juntos
ou impede combinações inválidas.

\ilustracaopendente{7-2-anatomia-store}{Diagrama da anatomia da store em três
faixas. A primeira mostra o contrato TypeScript; a segunda, o estado inicial
``isOpen false'' e ``initialEmail null''; a terceira, as ações ``openInvite'' e
``closeInvite'' chamando set. Setas devem mostrar que componentes enviam intenções
às ações e recebem somente as parcelas selecionadas.}

\section{Conectando os componentes}

O botão da página assina somente a ação:

\begin{lstlisting}
export function InviteButton() {
  const openInvite = useInviteStore((store) => store.openInvite)

  return (
    <button type="button" onClick={() => openInvite()}>
      Convidar pessoa
    </button>
  )
}
\end{lstlisting}

Uma sugestão da busca pode fornecer o endereço conhecido:

\begin{lstlisting}
openInvite('maria@empresa.com')
\end{lstlisting}

Já o componente visual seleciona os dados e a ação que utiliza:

\begin{lstlisting}
export function InviteDialog() {
  const isOpen = useInviteStore((store) => store.isOpen)
  const initialEmail = useInviteStore(
    (store) => store.initialEmail,
  )
  const closeInvite = useInviteStore(
    (store) => store.closeInvite,
  )

  if (!isOpen) return null

  return (
    <Dialog onClose={closeInvite}>
      <InviteForm initialEmail={initialEmail} />
    </Dialog>
  )
}
\end{lstlisting}

O modal pode ficar próximo da raiz visual para lidar corretamente com camada,
foco e acessibilidade. Os botões que o abrem não precisam ser ancestrais dele.
A store coordena o comportamento; ela não decide onde o HTML será renderizado.

\section{Acesso fora de componentes React}

O hook criado por \codigo{create} também expõe uma API da store. Em um evento do
navegador, podemos ler o estado atual sem chamar um hook:

\begin{lstlisting}[caption={Atalho global usando a API imperativa.}]
export function installInviteKeyboardShortcut() {
  function handleKeyDown(event: KeyboardEvent) {
    if (event.key !== 'Escape') return

    const store = useInviteStore.getState()
    if (store.isOpen) store.closeInvite()
  }

  document.addEventListener('keydown', handleKeyDown)

  return () => {
    document.removeEventListener('keydown', handleKeyDown)
  }
}
\end{lstlisting}

\codigo{getState()} devolve a fotografia atual da store; não cria uma assinatura
React. O retorno da função remove o listener e evita registrar o mesmo atalho
mais de uma vez durante remontagens ou recarregamentos de desenvolvimento.

Esse acesso deve ser pontual. Importar a store em qualquer arquivo é tecnicamente
possível, mas dependências imperativas demais escondem o fluxo. Eventos externos,
integrações e testes são usos claros; componentes devem preferir o hook com
selector.

\ilustracaopendente{7-3-sequencia-escape}{Sequência em quatro passos: modal
aberto; pessoa pressiona Escape; listener do documento consulta a store; ação
closeInvite altera isOpen para false e o componente deixa de renderizar o modal.
Destacar que o listener não manipula diretamente o DOM nem conhece o componente
visual.}

\section{O que ainda não pertence a esta store}

O texto digitado em cada campo pode ficar no formulário. O envio e a resposta do
servidor podem ser tratados como mutação remota. A visibilidade de outro modal
não deve entrar aqui apenas porque também é booleana.

Uma store chamada \codigo{useGlobalStore} tende a receber qualquer novo estado.
\codigo{useInviteStore} estabelece uma fronteira: coordenação do fluxo de
convite. O nome ajuda a equipe a dizer não a responsabilidades estranhas.

Também não precisamos introduzir persistência, middleware ou slices. Esses
recursos serão úteis quando um problema concreto os justificar. A primeira store
deve permanecer legível em uma única tela.

\section{Exercício: feche com Escape}

\begin{tarefa}
Implemente o fechamento do modal pela tecla Escape sem adicionar um listener
dentro de \codigo{InviteDialog} e sem manipular o DOM para esconder o elemento.
Instale o atalho uma única vez e garanta a remoção do listener.
\end{tarefa}

Depois, responda:

\begin{enumerate}
  \item por que \codigo{closeInvite} pertence ao contrato da store;
  \item qual diferença existe entre o hook e \codigo{getState()};
  \item onde o atalho deve ser instalado e removido;
  \item quais dados permanecem locais ao formulário;
  \item que nome genérico tornaria a fronteira menos clara.
\end{enumerate}

\espacoresposta{9}

\begin{resumo}
\begin{itemize}
  \item Um bom primeiro store começa por um contrato pequeno.
  \item Estado inicial e ações formam sua API pública.
  \item A forma \codigo{create<T>()(...)} mantém o contrato tipado.
  \item \codigo{set} combina atualizações no primeiro nível da store.
  \item Ações semânticas preservam regras e comunicam intenção.
  \item Componentes usam selectors; eventos externos podem usar
    \codigo{getState()} com moderação.
\end{itemize}
\end{resumo}

No próximo capítulo, os filtros do Dashboard mostrarão por que atualizar estado
não é apenas expor setters: as regras de mudança também precisam de um lugar.
